\begin{abstract}
\noindent Exoplanets and constructed worlds whose air and water resemble Earth's, while their gravity, size,
day or starlight differ, are studied with models built for Earth, in which an Earth value left in the code runs
without failing. Four such models were adapted to a breathable Moon, which changes gravity, column mass, scale
height and day length together: ExoPlaSim~3.4.2, CM1 release~22 with WRF-ELEC's electrification, SWAN~41.51
and SWASH~12.01. The adaptation found Earth values in four forms (constants, units of time, geometry, and fits
beyond their range) and five defects in the released codes. Three defects produce little or no error at
Earth's values and affect other users' runs: ExoPlaSim's Sun clock, which skipped back 32.5\textdegree\ each
lunar model day;\claim{H23} its convective-cover unit, which also biases synchronous runs;\claim{W06} and its
default orbit, which replaced the Moon's tilt with Earth's. General repairs correct units, apply scaling
laws or remove defects and serve other worlds unchanged; a method computes the curved geometry for each world's
column; and calibrations to an in-house line-by-line model hold near the design case. Identity, reproduction, similarity and benchmark checks test the code; the lunar climate is
unobserved. On the case, the general repairs moved the global mean by a few kelvin and the distribution of rain
further, the calibrations and PlaSim's Earth-tuned cloud and longwave cooling moved it most,\claim{G05} and the
storms' freezing level stood five to six times higher than Earth's, their lightning rarer and
larger.\claim{G09}\claim{G10} Building those parts from the reference's physics and from cloud-resolving runs
would narrow the Moon's uncertainty and widen the models' reach together.
\end{abstract}
